\documentclass[11pt]{article}

\usepackage[margin=1in]{geometry}
\usepackage[T1]{fontenc}
\usepackage[utf8]{inputenc}
\usepackage{lmodern}
\usepackage{amsmath}
\usepackage{amssymb}
\usepackage{booktabs}
\usepackage{array}
\usepackage{longtable}
\usepackage{enumitem}
\usepackage{xcolor}
\usepackage{hyperref}
\usepackage{CJKutf8}

\hypersetup{
  colorlinks=true,
  linkcolor=blue,
  urlcolor=blue,
  citecolor=blue
}

\setlist[itemize]{leftmargin=1.5em}
\setlist[enumerate]{leftmargin=1.8em}
\setlength{\parskip}{0.45em}
\setlength{\parindent}{0pt}

\title{MiniOneRec 后续研究\\进度日志}
\author{工作记录文档}
\date{最后更新：2026-04-10}

\begin{document}
\begin{CJK*}{UTF8}{gbsn}
\maketitle

\section{文档用途}

这份文档是当前 MiniOneRec 后续研究的持续更新日志，不是正式论文草稿。它的主要目标有四个：

\begin{itemize}
  \item 持续同步当前研究主线与最新实验结论；
  \item 记录每一轮实验实际做了什么、改了什么；
  \item 区分已经比较稳定的事实和仍然开放的想法；
  \item 为后续正式写论文保留一条清晰、可回溯的叙事链条。
\end{itemize}

\section{当前研究状态}

\subsection{当前方向}

当前研究方向可以概括为：

\begin{quote}
用图结构作为协同信息的载体，并以层级感知的方式，把 graph-structured collaborative information 融合进 MiniOneRec 的语义 SID 构建过程。
\end{quote}

这里的重点不是比较哪一种 graph encoder 更强，而是回答三个核心问题：

\begin{enumerate}
  \item 什么样的图适合承载协同信息？
  \item 图结构协同信息应该怎样在不同 SID level 上体现层级感知？
  \item 图协同信息应该如何进入 MiniOneRec 的纯语义 SID，而不破坏语义 backbone 的稳定性？
\end{enumerate}

\subsection{为什么转向这个方向}

当前方向主要来自以下几条逐步稳定下来的观察：

\begin{itemize}
  \item SID collision 存在，但看起来不是主矛盾；
  \item 当前 SID 主要由文本语义驱动，缺少直接协同约束；
  \item 大量错误表现为“prefix 对了，但 leaf 错了”，说明 local leaf ambiguity 更像核心瓶颈；
  \item 早期 naive front-end collaborative fusion 不但没有稳定提升，反而可能破坏量化结构。
\end{itemize}

\section{当前较稳定的事实}

到目前为止，以下结论已经足够稳定，可以继续指导方法设计：

\begin{enumerate}
  \item \textbf{局部歧义是真问题。} 模型在进入正确语义区域后，最后一层 leaf 仍然经常选错。
  \item \textbf{图协同信号是有用的。} 即使还没有改训练版 SID，graph-structured collaborative signal 已经能在 top-1 SID error 上提供额外判别力。
  \item \textbf{中尺度图视图很重要。} 最有价值的协同视图既不是纯 global，也不是纯 last-step local。
  \item \textbf{轻量净化是必要的。} raw collaborative graph 往往弱于 purified graph。
  \item \textbf{论文启发的频域中尺度图比第一版 heuristic middle-view 更强。}
  \item \textbf{当前仓库默认的 \texttt{sid-train} 超参数与上游 MiniOneRec 不一致。} 真正能复现低 collision SID 的，是 upstream-style 的 \texttt{RQ-VAE} 训练制度，而不是当前默认的 \texttt{epochs=10, batch\_size=256}。
  \item \textbf{不能只看 \texttt{sid-train} collision。} 最终真正有意义的是 \texttt{sid-generate} 后的 \texttt{index.json} collision；train-stage 排名和 final index 排名可能并不一致。
\end{enumerate}

\section{被削弱或阶段性放下的方向}

\subsection{Naive 前端协同融合}

之前的实验说明，粗糙的前端语义加协同融合会破坏 SID 结构并伤害量化稳定性。这并不意味着所有前端协同融合都是错的，只是否掉了那一版过于粗糙的融合方式。

\subsection{把后验局部修正当成最终论文主线}

局部后验修正仍然是很有用的诊断工具和原型验证工具，但单独作为最终论文主方法太弱。更有价值的主线，还是把图结构协同信息以层级感知方式写回 SID 构建本身。

\section{实验时间线}

\subsection{阶段 0：问题收敛}

这一阶段的主要贡献还是概念层面的。项目从一个较宽泛的“往 SID 里加协同信息”方向，逐渐收敛到更明确的表述：

\begin{quote}
图结构承载的协同信息，应该以层级感知且保结构的方式进入 SID 构建。
\end{quote}

这一阶段同时也固定了三个核心问题，并明确了图在当前研究中的身份：

\begin{quote}
图不是主要比较对象，而是 collaborative information 的结构化载体。
\end{quote}

\subsection{阶段 1：初始 Graph-Bank Probe（2026-04-09）}

\subsubsection*{做了什么}

我们新增了一套独立 probe，不修改 MiniOneRec 原有训练或评测主链，只在已有 top-1 SID 错误样本上测试不同 graph view 的判别能力。测试的视图包括：

\begin{itemize}
  \item coarse collaborative graph；
  \item local transition graph；
  \item heuristic middle-resolution views：
  \begin{itemize}
    \item diffusion residual；
    \item band-pass proxy；
    \item community-aware view。
  \end{itemize}
\end{itemize}

其中适用时分别比较 raw 与 purified 版本。

\subsubsection*{主要结论}

第一轮里最有价值的 middle-resolution views 来自 purified heuristic views：

\begin{itemize}
  \item Industrial：最好的 middle view 是 \texttt{mid\_band\_pass\_purified}，\texttt{same\_l2 = 0.2747}
  \item Office：最好的 middle view 是 \texttt{mid\_diffusion\_purified}，\texttt{same\_l2 = 0.2381}
\end{itemize}

这一轮的意义是明确了：

\begin{quote}
middle-resolution collaborative structure 值得被当成一个一等设计对象，而不是顺手派生出来的附属信号。
\end{quote}

\subsubsection*{说明了什么}

这轮结果支持下面这个解释：

\begin{quote}
对深层 leaf ambiguity 最有用的协同信息，既不是纯 global，也不是纯 local。中尺度视图是必要的。
\end{quote}

\subsection{阶段 2：论文模块移植 Probe（2026-04-09）}

\subsubsection*{做了什么}

第二轮 probe 依然保持为独立实验模块，但移植了两类更接近相关工作的思想：

\begin{itemize}
  \item \textbf{PRISM 风格的 semantic-anchor purification}
  \item \textbf{GSPRec / FaGSP 风格的频域 middle-resolution graph views}
\end{itemize}

这一轮仍然没有改 SID 训练本体，而是在同样的 top-1 error probe 框架下测试更强的 graph carrier。

\subsubsection*{主要结果}

\begin{longtable}{p{0.26\linewidth}p{0.16\linewidth}p{0.16\linewidth}p{0.24\linewidth}}
\toprule
视图类型 & Industrial 最好 \texttt{same\_l2} & Office 最好 \texttt{same\_l2} & 说明 \\
\midrule
第一版 heuristic middle-view & 0.2747 & 0.2381 & 上一轮最强 purified proxy \\
移植后的 spectral middle-view & 0.5401 & 0.7619 & 都来自 \texttt{fagsp\_mid\_base} \\
\bottomrule
\end{longtable}

更具体地说：

\begin{itemize}
  \item Industrial：
  \begin{itemize}
    \item 上一轮最强中尺度视图：\texttt{mid\_band\_pass\_purified = 0.2747}
    \item 当前最强移植视图：\texttt{fagsp\_mid\_base = 0.5401}
    \item 第二强移植视图：\texttt{gsprec\_mid\_prism = 0.5247}
  \end{itemize}
  \item Office：
  \begin{itemize}
    \item 上一轮最强中尺度视图：\texttt{mid\_diffusion\_purified = 0.2381}
    \item 当前最强移植视图：\texttt{fagsp\_mid\_base = 0.7619}
    \item 第二强移植视图：\texttt{gsprec\_mid\_prism = 0.6429}
  \end{itemize}
\end{itemize}

\subsubsection*{说明了什么}

这一轮是目前为止最强的一次正信号：

\begin{quote}
当前的 \texttt{G\_mid} 很可能应该被设计为一个真正的 graph-spectral middle-resolution collaborative view，而不是只依赖 heuristic diffusion / band-pass proxy。
\end{quote}

同时，这轮结果也说明：

\begin{quote}
semantic-anchor purification 本身是有帮助的，但真正的大幅增益主要来自更强的中尺度结构表达，而不只是 anchor-only denoising。
\end{quote}

\subsection{阶段 3：训练脚手架校准（过渡阶段）}

在进入正式训练验证之前，我们做过几轮训练脚手架和配置校准，用来确认两件事：

\begin{itemize}
  \item 独立实验分支可以在不破坏 baseline 主链的前提下稳定运行；
  \item 真正有意义的方法比较必须建立在与 MiniOneRec 上游一致的训练制度上。
\end{itemize}

这一阶段的详细过程不再作为当前主线证据保留。对当前研究真正有价值的结论只有一条：

\begin{quote}
在做 graph-aware SID 对比之前，必须先解决 baseline provenance 与训练制度对齐问题。
\end{quote}

\subsection{阶段 4：SID provenance 追踪与上游风格复现（2026-04-09）}

\subsubsection*{做了什么}

为了确认当前正式 Industrial SID 的真实来源，我们专门做了两件事：

\begin{itemize}
  \item 对比本地迁移版与上游 MiniOneRec 官方仓库中的 \texttt{rqvae.py}、\texttt{trainer.py}、\texttt{generate\_indices.py} 核心逻辑；
  \item 在本地主线代码上，按上游官方训练超参数重新跑一遍 Industrial 的 \texttt{RQ-VAE -> sid-generate}。
\end{itemize}

\subsubsection*{主要结果}

对比上游后发现：

\begin{itemize}
  \item 核心 \texttt{sid-train / sid-generate} 代码逻辑没有被迁移改坏；
  \item 真正的关键差异在于训练制度：
  \begin{itemize}
    \item 上游官方：\texttt{epochs = 10000, batch\_size = 20480}
    \item 当前本地默认：\texttt{epochs = 10, batch\_size = 256}
  \end{itemize}
\end{itemize}

随后，按上游风格重跑 Industrial：

\begin{itemize}
  \item \texttt{sid-train} 最佳 collision 降到 \(\approx 0.0993\)；
  \item 继续跑 \texttt{sid-generate} 后，最终 \texttt{index.json} collision 变为 \(\approx 0.0043407488\)；
  \item 这个值与仓库中当前正式 Industrial index 的 collision 完全一致。
\end{itemize}

\subsubsection*{说明了什么}

这一轮非常关键，因为它修正了我们对基线系统的理解：

\begin{quote}
当前正式 Industrial SID 是可以被当前主线代码复现出来的，问题不在于迁移时改坏了主逻辑，而主要在于本地默认训练超参数偏离了上游 MiniOneRec。
\end{quote}

这也意味着，后续所有 graph-aware 实验如果还想保持说服力，就必须在 \textbf{upstream-aligned} 训练制度下进行。

\subsection{阶段 6：Upstream-aligned MGR-SID v1 训练与生成（2026-04-09 / 2026-04-10）}

\subsubsection*{做了什么}

在确认基线 provenance 后，我们把 \texttt{MGR-SID v1} 的训练实验也切换到了与上游对齐的训练制度：

\begin{itemize}
  \item \texttt{epochs = 10000}
  \item \texttt{batch\_size = 20480}
  \item Industrial only
  \item 三组对照：
  \begin{itemize}
    \item baseline
    \item uniform graph regularization
    \item hierarchy-aware graph regularization
  \end{itemize}
\end{itemize}

训练结束后，我们没有停留在 \texttt{sid-train} 指标，而是对每个 best checkpoint 单独跑了实验版 generate，比较最终 \texttt{index.json} collision。

\subsubsection*{主要结果}

\paragraph{Generate-stage final collision}

\begin{itemize}
  \item baseline：\(\approx 0.0035268584\)
  \item uniform\_reg：\(\approx 0.0059685296\)
  \item hierarchy\_reg：\(\approx 0.0032555616\)
\end{itemize}

其中最关键的是：

\begin{itemize}
  \item \texttt{hierarchy\_reg} 成为三者里最终最好的 SID；
  \item baseline 也比当前仓库里原有正式 Industrial index 的 \(\approx 0.0043407488\) 更低；
  \item 说明这轮 upstream-aligned 训练本身已经把语义 baseline 推到了更健康的 regime。
\end{itemize}

\subsubsection*{说明了什么}

这是当前研究里最重要的新转折之一：

\begin{quote}
当前最重要的判断标准应该是 final generated SID，而不是单独盯住 raw \texttt{sid-train} collision；在这个更合理的判据下，当前这版 hierarchy-aware graph regularization 已经给出了第一轮清晰正信号。
\end{quote}

更具体地说，这一轮至少说明了三件事：

\begin{itemize}
  \item \texttt{hierarchy-aware} 方向不是空想，它已经在最终 Industrial SID 上给出正结果；
  \item 后续评估应该优先围绕 \textbf{final generated SID} 展开，而不是再把全部判断压在 train-stage collision 上；
  \item 当前实验主线已经具备继续推进的依据，不再只是 probe 级别的希望信号。
\end{itemize}

\section{当前证据的整体解释}

把目前的证据串起来，可以得到一个更清楚的图景：

\begin{enumerate}
  \item 当前 semantic SID 足以提供较强的 coarse routing；
  \item 剩余的主要困难集中在 local leaf-level discrimination；
  \item graph-structured collaborative signal 对这个瓶颈是有帮助的；
  \item 当前最强候选是 \textbf{spectral middle-resolution collaborative graph view}；
  \item 当前仓库默认 \texttt{sid-train} 配置不能代表正式 MiniOneRec baseline，所有方法对比必须切到 upstream-aligned regime；
  \item 在 upstream-aligned regime 下，\texttt{hierarchy-aware} 已经在 final generated SID 上给出了第一轮正信号；
  \item 这说明图信号本身有用，而且图信号如何与 \texttt{sid-generate} 的 collision resolution 交互，是方法设计里非常关键的一环。
\end{enumerate}

\section{当前最佳工作假设}

当前最可信的工作假设是：

\begin{quote}
MiniOneRec 应该继续保留 semantic SID 作为 backbone，同时加入 hierarchy-aware graph regularization；其中 middle-resolution spectral collaborative graph 是降低 local leaf ambiguity 的关键额外信号，而真正的有效性要通过 \texttt{sid-generate} 后的 final SID 来判断。
\end{quote}

目前最强的 \texttt{G\_mid v1} 候选是：

\begin{itemize}
  \item 主候选：\texttt{fagsp\_mid\_base}
  \item 备选：\texttt{gsprec\_mid\_prism}
\end{itemize}

当前最合理的控制组地位也更清楚了：

\begin{itemize}
  \item \texttt{hierarchy\_reg}：当前主方法候选
  \item \texttt{baseline}：强语义基线
\end{itemize}

\section{当前开放问题}

即使有了新一轮 probe，下面这些问题仍然没有完全解决：

\begin{enumerate}
  \item 当前 final-SID 改善是否能在第二个 seed 上保持？
  \item baseline 与 hierarchy\_reg 的最终 index 差异，是否真的主要集中在 \texttt{same\_l2} / local ambiguity bucket？
  \item 当前 \texttt{hierarchy\_reg} 的收益究竟来自更好的 leaf disambiguation，还是更好的 collision repair compatibility？
  \item semantic anchoring 在训练时应该如何使用，才能避免把 graph signal 限制得过死？
  \item graph regularization 最佳作用位置到底是 encoder latent、quantized residual，还是 cumulative quantized representation？
  \item 当前 final-SID 的改进，最终能否转化到 HR@1、NDCG@10 等主推荐指标上？
  \item 什么时候适合把 published baselines 拉进主表，避免在机制还没站稳前过早铺开大规模复现？
\end{enumerate}

\section{下一步行动}

当前最合理的 next actions 是：

\begin{enumerate}
  \item 保留当前 graph bank，不再频繁改 carrier：
  \begin{itemize}
    \item \texttt{G\_coarse}：purified coarse collaborative graph；
    \item \texttt{G\_mid}：\texttt{fagsp\_mid\_base}；
    \item \texttt{G\_local}：purified transition graph。
  \end{itemize}
  \item 优先围绕 final generated SID 做分析，而不是再把注意力只放在 \texttt{sid-train} collision 上；
  \item 下一轮最优先比较：
  \begin{itemize}
    \item baseline；
    \item hierarchy-aware graph regularization。
  \end{itemize}
  \item 先补两类低成本但高信息量分析：
  \begin{itemize}
    \item baseline vs hierarchy 的 final index bucket analysis；
    \item baseline vs hierarchy 的 item-level changed cases。
  \end{itemize}
  \item 如果第二个 seed 也保持正向，再考虑做：
  \begin{itemize}
    \item 更细的 integration 改写；
    \item downstream recommendation evaluation；
    \item published baselines 主表准备。
  \end{itemize}
\end{enumerate}

\section{当前关键材料位置}

当前阶段最关键的材料包括：

\begin{itemize}
  \item 任务对齐文档：\path{idea-discovery/.../CURRENT_TASK_ALIGNMENT.md}
  \item 第一轮 probe 记录：\path{idea-discovery/.../13_initial_probe_run_2026-04-09.md}
  \item 论文模块移植 probe 记录：\path{idea-discovery/.../14_paper_transplant_probe_run_2026-04-09.md}
  \item 上游风格 SID 复现记录：\path{research-progress-log/experiment_launches/2026-04-09_upstream_style_rqvae_result.md}
  \item upstream-aligned MGR-SID v1 总记录：\path{research-progress-log/experiment_launches/2026-04-09_mgr_sid_v1_upstream/README.md}
  \item upstream-aligned MGR-SID v1 结果记录：\path{research-progress-log/experiment_launches/2026-04-09_mgr_sid_v1_upstream/RESULTS.md}
  \item 第一轮 probe 脚本：\path{scripts/experiment_mgr_sid_graph_bank_probe.py}
  \item 模块移植 probe 脚本：\path{scripts/experiment_mgr_sid_paper_transplant_probe.py}
  \item 实验版训练脚本：\path{scripts/experiment_mgr_sid_v1_train.py}
  \item 实验版生成脚本：\path{scripts/experiment_mgr_sid_v1_generate.py}
\end{itemize}

\end{CJK*}
\end{document}
